\section{Sujet n\textsuperscript{o}\,19}
\noindent\begin{minipage}{0.5\linewidth}\footnotesize Office du Baccalauréat du Cameroun\end{minipage}%
\hfill\begin{minipage}{0.45\linewidth}\footnotesize\raggedleft Examen : \textbf{Baccalauréat} · Série : \textbf{C/E}\\
Épreuve : \textbf{Mathématiques} · Durée : \textbf{4\,h} · Coef. : \textbf{7}\end{minipage}
\par\smallskip{\color{onbo}\rule{\linewidth}{1pt}}

\subsection*{Partie A — Évaluation des ressources \hfill (15 points)}

\noindent\textbf{Exercice 1.} \hfill\textit{(3 points)}\\
À l'hôpital régional de Garoua, le service des urgences compte $3$ lits de
réanimation et $5$ lits ordinaires (soit $8$ lits au total). On choisit au hasard
$3$ lits pour une visite d'inspection. Soit $X$ la variable aléatoire égale au
nombre de lits de réanimation parmi les $3$ lits inspectés.
\begin{enumerate}[label=\arabic*)]
\item Déterminer la loi de probabilité de $X$. \hfill\textit{(1,5 pt)}
\item Calculer l'espérance mathématique $E(X)$. \hfill\textit{(0,75 pt)}
\item Calculer la variance $V(X)$. \hfill\textit{(0,75 pt)}
\end{enumerate}

\smallskip
\noindent\textbf{Exercice 2.} \hfill\textit{(3 points)}\\
Dans un service, on modélise le nombre de patients hospitalisés le matin du jour
$n$ par une suite $(u_n)$. Le premier jour, $u_0=80$ patients. Chaque jour,
\SI{20}{\percent} des patients sortent (guéris) et $30$ nouveaux patients sont
admis. On a donc $u_{n+1}=0{,}8\,u_n+30$.
\begin{enumerate}[label=\arabic*)]
\item Calculer $u_1$ et $u_2$. \hfill\textit{(0,5 pt)}
\item On pose $v_n=u_n-150$. Montrer que $(v_n)$ est une suite géométrique dont on
précisera la raison et le premier terme. \hfill\textit{(1 pt)}
\item Exprimer $v_n$ puis $u_n$ en fonction de $n$. \hfill\textit{(1 pt)}
\item Déterminer la limite de $(u_n)$ et l'interpréter. \hfill\textit{(0,5 pt)}
\end{enumerate}

\smallskip
\noindent\textbf{Exercice 3.} \hfill\textit{(4 points)}\\
La concentration (en mg/L) d'un médicament dans le sang d'un patient est modélisée
par une fonction $y$ du temps $t\ge0$ (en heures), solution de l'équation
différentielle $(E)\ :\ y'+2y=0$.
\begin{enumerate}[label=\arabic*)]
\item Résoudre l'équation différentielle $(E)$. \hfill\textit{(1 pt)}
\item On injecte au temps $t=0$ une dose telle que $y(0)=10$. Déterminer la
fonction $f$ correspondante. \hfill\textit{(0,75 pt)}
\item Étudier les variations de $f$ sur $[0\,;\,+\infty[$ et déterminer
$\lim_{t\to+\infty}f(t)$. \hfill\textit{(1,25 pt)}
\item Le médicament n'est plus actif lorsque sa concentration devient inférieure à
$\SI{1}{mg/L}$. Déterminer, à l'heure près par excès, le temps au bout duquel le
médicament cesse d'être actif. \hfill\textit{(1 pt)}
\end{enumerate}

\smallskip
\noindent\textbf{Exercice 4.} \hfill\textit{(5 points)}\\
L'espace est muni d'un repère orthonormé direct $(O\,;\,\vec\imath,\vec\jmath,
\vec k)$. On considère les points $A(2\,;\,0\,;\,0)$, $B(0\,;\,3\,;\,0)$,
$C(0\,;\,0\,;\,4)$ et $D(2\,;\,3\,;\,4)$.
\begin{enumerate}[label=\arabic*)]
\item Déterminer les coordonnées du produit vectoriel
$\overrightarrow{AB}\wedge\overrightarrow{AC}$. \hfill\textit{(1,25 pt)}
\item En déduire que le plan $(ABC)$ a pour équation cartésienne
$6x+4y+3z-12=0$. \hfill\textit{(1 pt)}
\item Calculer la distance du point $D$ au plan $(ABC)$. \hfill\textit{(1,25 pt)}
\item Calculer l'aire du triangle $ABC$, puis en déduire le volume du tétraèdre
$ABCD$. \hfill\textit{(1,5 pt)}
\end{enumerate}

\subsection*{Partie B — Évaluation des compétences \hfill (5 points)}

\noindent\textit{Le directeur de l'hôpital de Garoua te sollicite, en tant
qu'élève de Terminale~C, pour optimiser la gestion. Le service de pédiatrie veut
construire une réserve de stockage à \emph{base carrée} et toit ouvert, de volume
$\SI{32}{\metre\cubed}$, en minimisant la surface de tôle (le fond carré et les
quatre parois verticales). Par ailleurs, le stock de poches de sang, de
$2000$ unités aujourd'hui, diminue de \SI{15}{\percent} chaque semaine ; on veut
savoir quand il passera sous le seuil de $500$ unités. Enfin, sur $10$ patients
admis aux urgences, chacun présente, indépendamment, une probabilité $0{,}2$
d'avoir besoin d'une transfusion.}

\begin{enumerate}[label=\textbf{Tâche \arabic*.},leftmargin=2.4em]
\item Déterminer les dimensions de la réserve qui minimisent la surface de tôle et
préciser cette surface minimale. \hfill\textit{(1,5 pt)}
\item Déterminer le nombre de semaines au bout duquel le stock de poches de sang
passera sous le seuil de $500$ unités. \hfill\textit{(1,5 pt)}
\item Déterminer la probabilité qu'\emph{exactement} $3$ des $10$ patients admis
aient besoin d'une transfusion. \hfill\textit{(1,5 pt)}
\end{enumerate}
\par\smallskip{\footnotesize\itshape Présentation générale : 0,5 point.}

% ════════════════════════════════════════════════════
\corrige{du sujet 19}

\noindent\textbf{Partie A — Exercice 1.}
1) $\binom{8}{3}=56$ choix. $X\in\{0,1,2,3\}$ et
$P(X{=}k)=\dfrac{\binom3k\binom5{3-k}}{56}$ :
$P(0)=\frac{10}{56}=\frac5{28}$, $P(1)=\frac{30}{56}=\frac{15}{28}$,
$P(2)=\frac{15}{56}$, $P(3)=\frac{1}{56}$. (Vérif. : $\frac{20+30+15+1}{56}=1$.)
2) $E(X)=0\cdot\frac{10}{56}+1\cdot\frac{30}{56}+2\cdot\frac{15}{56}+3\cdot
\frac1{56}=\frac{30+30+3}{56}=\frac{63}{56}=\frac98=1{,}125$.
3) $E(X^2)=0+1\cdot\frac{30}{56}+4\cdot\frac{15}{56}+9\cdot\frac1{56}
=\frac{30+60+9}{56}=\frac{99}{56}$.
$V(X)=E(X^2)-E(X)^2=\frac{99}{56}-\frac{81}{64}$. Au même dénominateur $448$ :
$\frac{792}{448}-\frac{567}{448}=\frac{225}{448}\approx0{,}502$.

\noindent\textbf{Exercice 2.}
1) $u_1=0{,}8\times80+30=64+30=94$ ; $u_2=0{,}8\times94+30=75{,}2+30=105{,}2$.
2) $v_n=u_n-150$ : $v_{n+1}=u_{n+1}-150=0{,}8u_n+30-150=0{,}8u_n-120
=0{,}8(u_n-150)=0{,}8\,v_n$. Donc $(v_n)$ est géométrique de raison $q=0{,}8$ et de
premier terme $v_0=u_0-150=80-150=-70$.
3) $v_n=-70\times0{,}8^{\,n}$, donc $u_n=v_n+150=150-70\times0{,}8^{\,n}$.
4) Comme $0<0{,}8<1$, $0{,}8^{\,n}\to0$, donc $u_n\to150$. À long terme, le service
se stabilise autour de $150$ patients hospitalisés.

\noindent\textbf{Exercice 3.}
1) $(E)\ :\ y'=-2y$ ; solutions $y(t)=K\mathrm{e}^{-2t}$, $K\in\R$.
2) $f(0)=K=10$, donc $f(t)=10\,\mathrm{e}^{-2t}$.
3) $f'(t)=-20\,\mathrm{e}^{-2t}<0$ : $f$ est strictement décroissante sur
$[0\,;\,+\infty[$, de $f(0)=10$ vers $\lim_{t\to+\infty}f(t)=0$ (la droite $y=0$
est asymptote horizontale).
4) On cherche $t$ tel que $f(t)<1$ : $10\mathrm{e}^{-2t}<1\Leftrightarrow
\mathrm{e}^{-2t}<0{,}1\Leftrightarrow-2t<\ln0{,}1\Leftrightarrow
t>\frac{\ln10}{2}\approx\frac{2{,}3026}{2}\approx1{,}15$. Le médicament cesse
d'être actif au bout de \textbf{$2$ heures} (premier entier d'heures dépassant
$1{,}15$).

\noindent\textbf{Exercice 4.}
1) $\overrightarrow{AB}(-2\,;\,3\,;\,0)$, $\overrightarrow{AC}(-2\,;\,0\,;\,4)$.
$\overrightarrow{AB}\wedge\overrightarrow{AC}=
\big(3\cdot4-0\cdot0\,;\,0\cdot(-2)-(-2)\cdot4\,;\,(-2)\cdot0-3\cdot(-2)\big)
=(12\,;\,8\,;\,6)$.
2) Le vecteur normal $(12\,;\,8\,;\,6)$, ou son colinéaire $(6\,;\,4\,;\,3)$, donne
$(ABC)\ :\ 6x+4y+3z+d=0$. Avec $A(2\,;\,0\,;\,0)$ : $12+d=0$, $d=-12$. D'où
$6x+4y+3z-12=0$.
3) $d(D,(ABC))=\dfrac{|6\cdot2+4\cdot3+3\cdot4-12|}{\sqrt{6^2+4^2+3^2}}
=\dfrac{|12+12+12-12|}{\sqrt{36+16+9}}=\dfrac{24}{\sqrt{61}}=\dfrac{24\sqrt{61}}{61}
\approx3{,}07$.
4) Aire de $ABC$ : $\mathcal{A}=\frac12\|\overrightarrow{AB}\wedge
\overrightarrow{AC}\|=\frac12\sqrt{12^2+8^2+6^2}=\frac12\sqrt{144+64+36}
=\frac12\sqrt{244}=\sqrt{61}$.
Volume du tétraèdre : $V=\frac13\times\mathcal{A}\times d(D,(ABC))
=\frac13\times\sqrt{61}\times\frac{24}{\sqrt{61}}=\frac{24}{3}=8$. Donc
$V=\SI{8}{}$ unités de volume.

\smallskip
\noindent\textbf{Partie B — Situation.}
\emph{Tâche 1.} Base carrée de côté $x$, hauteur $h$ ; volume $x^2h=32$, donc
$h=\frac{32}{x^2}$. Surface de tôle (fond + $4$ parois, pas de toit) :
$S(x)=x^2+4xh=x^2+4x\cdot\frac{32}{x^2}=x^2+\frac{128}{x}$.
$S'(x)=2x-\frac{128}{x^2}=\frac{2x^3-128}{x^2}=0\Leftrightarrow x^3=64
\Leftrightarrow x=4$. $S'$ change de $-$ à $+$ en $4$ : minimum. Alors
$h=\frac{32}{16}=2$. Dimensions : base $4\times4$, hauteur $2$ ($x=\SI{4}{\metre}$,
$h=\SI{2}{\metre}$) ; surface minimale $S(4)=16+\frac{128}{4}=16+32
=\SI{48}{\metre\squared}$.

\emph{Tâche 2.} Stock de la semaine $n$ : $w_n=2000\times0{,}85^{\,n}$. On résout
$w_n<500$, soit $0{,}85^{\,n}<\frac14$, d'où $n>\frac{\ln(1/4)}{\ln0{,}85}
=\frac{-\ln4}{\ln0{,}85}\approx\frac{-1{,}3863}{-0{,}1625}\approx8{,}53$. Le stock
passe sous $500$ unités à partir de la \textbf{$9^{\text{e}}$ semaine}
($w_9\approx2000\times0{,}2316\approx463$ unités).

\emph{Tâche 3.} Le nombre $Y$ de patients nécessitant une transfusion suit la loi
binomiale $\mathcal{B}(10\,;\,0{,}2)$. $P(Y=3)=\binom{10}{3}(0{,}2)^3(0{,}8)^7
=120\times0{,}008\times0{,}2097152\approx\mathbf{0{,}201}$, soit environ
\SI{20,1}{\percent}.
